\section{Thuật toán DPLL, suy luận phản chứng và kiểm tra nghiệm duy nhất}

\subsection{Thủ tục tìm mô hình DPLL}

Bộ giải trong \texttt{sat.py} được nhóm tự cài đặt theo thuật toán DPLL \cite{davis1962machine}. Trước khi tìm kiếm, hàm \texttt{solve()} khử literal lặp trong từng mệnh đề, bỏ mệnh đề hằng đúng và ghép các giả định (\texttt{assumptions}) thành mệnh đề đơn. Mỗi nhánh đệ quy dùng bản sao phép gán riêng để cô lập khi quay lui.

\begin{algorithm}[H]
\small\singlespacing
\caption{DPLL-Solve với hai chiến lược chọn biến}
\begin{algorithmic}[1]
\Require Công thức CNF $F$, tổng số biến $m$, phép gán $A$ (khởi tạo rỗng), chiến lược $h$
\Loop
\State Kiểm tra cờ hủy và hạn thời gian; nếu $F$ có mệnh đề rỗng thì trả UNSAT.
\State Thu thập tập $U$ các literal trong mệnh đề đơn; nếu $U=\varnothing$ thì thoát vòng lặp.
\For{mỗi $\ell\in U$}
\State Nếu $\ell$ ngược giá trị trong $A$ thì trả UNSAT; nếu biến chưa gán thì ghi vào $A$ và tăng số lan truyền.
\EndFor
\State Rút gọn $F$ theo $A$: bỏ mệnh đề đã thỏa và xóa literal sai khỏi mệnh đề còn lại.
\EndLoop
\If{$F$ rỗng}
\State Điền False cho mọi biến chưa gán trong $1,\ldots,m$; trả SAT kèm mô hình.
\EndIf
\State Đếm tần suất xuất hiện (gộp cả hai dấu) của từng biến còn lại trong $F$.
\State Với \texttt{ordered}, chọn chỉ số nhỏ nhất; với \texttt{frequency}, chọn tần suất lớn nhất (hòa chọn chỉ số nhỏ nhất). Gọi biến được chọn là $v$.
\State Tăng bộ đếm quyết định thêm 1.
\For{$b$ lần lượt nhận True, False}
\State Tạo $A' = A \cup \{v \mapsto b\}$, rút gọn $F$ theo $A'$ và gọi đệ quy DPLL-Solve.
\State Nếu nhánh con trả SAT thì trả mô hình; nếu thất bại thì tăng bộ đếm quay lui.
\EndFor
\State Trả UNSAT.
\end{algorithmic}
\end{algorithm}

Hai mệnh đề đơn đối dấu gây xung đột ngay trong vòng lan truyền. Khi $F$ rỗng, điền \texttt{False} cho các biến tự do vẫn bảo toàn nghiệm. Phép gán phân nhánh được ghi thẳng vào $A'$ nên \texttt{propagations} chỉ đếm đúng các phép gán bị ép buộc, còn thời gian được đo riêng cho từng lời gọi SAT.

\clearpage

\subsection{Kiểm chứng phán quyết bằng phản chứng và lưu mô hình nhân chứng}

Trước hết \texttt{LogicAgent} kiểm tra $KB_t$ khả thỏa để loại trường hợp cơ sở tri thức mâu thuẫn. Khi đó, $KB_t\land\neg C_i$ vô nghiệm chứng minh ô $i$ là tội phạm, còn $KB_t\land C_i$ vô nghiệm chứng minh ô $i$ vô tội. Nếu cả hai đều khả thỏa, mỗi mô hình tìm được đóng vai trò nhân chứng (\textit{witness}) cho một khả năng và ô $i$ vẫn ở trạng thái \texttt{UNKNOWN}.

\begin{algorithm}[H]
\small\singlespacing
\caption{Phân loại các ô chưa mở và kiểm chứng đề nghị}
\begin{algorithmic}[1]
\Require Trạng thái công khai \texttt{PublicState}; chế độ phân loại hoặc kiểm chứng literal đề nghị $p$
\State Dựng $KB_t$ và giải một lần; nếu UNSAT thì trả \texttt{INCONSISTENT}.
\If{chế độ \texttt{verify}}
\State Nếu $KB_t\land\neg p$ là UNSAT thì trả \texttt{ACCEPTED}.
\State Nếu $KB_t\land p$ là UNSAT thì trả \texttt{CONTRADICTED}; ngược lại trả \texttt{NOT\_PROVABLE}.
\EndIf
\State Từ mô hình cơ sở, đánh dấu giá trị (\texttt{True}/\texttt{False}) đã xuất hiện của từng biến chính.
\For{mỗi ô chưa mở theo thứ tự hàng}
\State Nếu $C_i$ đã xuất hiện đủ cả hai giá trị \texttt{True} và \texttt{False}, xuất \texttt{UNKNOWN} và sang ô tiếp theo.
\State Gọi $\ell$ là chiều giá trị chưa có mô hình nhân chứng; giải $KB_t\land\ell$.
\If{kết quả là UNSAT}
\State Xuất \texttt{CRIMINAL} nếu $\ell=\neg C_i$, ngược lại xuất \texttt{INNOCENT}.
\Else
\State Cập nhật giá trị của mọi biến chính từ mô hình mới và xuất \texttt{UNKNOWN}.
\EndIf
\EndFor
\end{algorithmic}
\end{algorithm}

Trong \texttt{LogicAgent.\_verdicts}, mô hình cơ sở ban đầu đã cho sẵn một giá trị thỏa mãn tại mọi biến chính, nên mỗi ô chưa mở chỉ cần tối đa một truy vấn SAT bổ sung (hoặc bằng 0 nếu các mô hình trước đã phủ đủ hai chiều). Bộ nhớ nhân chứng này chỉ tồn tại trong phạm vi một lượt duyệt $KB_t$, còn \texttt{verify()} kiểm tra độc lập từng đề nghị gửi lên nên có thể cần cả hai truy vấn giả định.

Với ví dụ $KB=(C_A\lor C_B)\land(\neg C_A\lor C_B)$, giả định $\neg C_B$ kéo theo cả $C_A$ lẫn $\neg C_A$ nên $B$ bắt buộc là tội phạm; trong khi $A$ nhận được cả \texttt{True} lẫn \texttt{False} khi $C_B=\mathrm{True}$, do đó trạng thái của $A$ là \texttt{UNKNOWN} và thẻ $A$ giữ nguyên úp.

\clearpage

\subsection{Vòng lặp mở thẻ, nhật ký suy luận và kiểm tra nghiệm duy nhất}

Hàm \texttt{next\_verdict()} dừng tại ô bắt buộc đầu tiên: \texttt{hint()} trả về ô gợi ý, còn \texttt{step()} gọi \texttt{submit()} để kiểm chứng lại, lật thẻ, ghi nhật ký (số bước, ô, phán quyết, danh sách ID lời khai đang mở, truy vấn SAT và lời khai mới) rồi dùng trạng thái công khai đã cập nhật cho bước kế tiếp. Khi mọi ô chưa mở đều còn \texttt{UNKNOWN}, vòng tự giải dừng mà không đoán mò.

Nhật ký thực tế của \texttt{case-3x3-01} (\texttt{ordered}) minh họa rõ quy trình này: ở bước 1 với hai lời khai B1 và C1, $KB_1$ gồm 9 biến chính, 39 biến phụ và 123 mệnh đề. Truy vấn cơ sở trả SAT, còn giả định $-7$ trả UNSAT sau 46 phép lan truyền và 0 quyết định, qua đó chứng minh A3 là tội phạm và mở lời khai A3 cho bước 2.

Để kiểm định bộ dữ liệu ngoại tuyến, \texttt{check\_uniqueness()} mã hóa toàn bộ $N^2$ lời khai thành $F$. Nếu $F$ vô nghiệm thì trả \texttt{INCONSISTENT}; nếu có mô hình $M$, hàm thêm mệnh đề chặn chỉ trên tập biến chính $P=\texttt{kb.primary.values()}$:

\[
B_M=\bigvee_{v\in P}\begin{cases}\neg v,&M(v)=\mathrm{True},\\v,&M(v)=\mathrm{False}.\end{cases}
\]

Mọi nghiệm của $F\land B_M$ buộc phải khác $M$ ở ít nhất một nhân vật: kết quả UNSAT xác nhận \texttt{UNIQUE}, còn SAT báo đa nghiệm (\texttt{MULTIPLE}). Chỉ chặn biến chính giúp tránh đếm nhầm khác biệt ở biến phụ Tseitin. Khả năng giải tuần tự từ tập mở đầu vẫn được kiểm tra riêng qua vòng suy luận công khai.

\subsection{Bộ kiểm thử tự động}

Lệnh \texttt{python -B -m unittest discover -s tests -v} chạy đạt toàn bộ 87/87 bài kiểm thử trên 13 mô-đun. Nhóm kiểm thử lõi đối chiếu vét cạn \texttt{evaluate()} với CNF, kiểm tra giả định, quay lui, điều kiện lồng, vùng rỗng và đăng ký luật mới; các mô-đun còn lại kiểm tra tính bất biến của \texttt{PublicState}, từ chối nước đoán mò, hủy worker nền, bộ thực nghiệm và giao diện.
